% ==============================================================================
% 07_criteria4_assessment.tex — Criteria 4: Student Assessment
% Score: 2
% ==============================================================================

\criterion{4}{การประเมินผู้เรียน}{Student Assessment}

% ------------------------------------------------------------------------------
\criterionsub{4.1 A variety of assessment methods are shown to be used and are shown to be constructively aligned to achieving the expected learning outcomes and the teaching and learning objectives.}

หลักสูตรกำหนดให้ทุกรายวิชาใช้วิธีการประเมินที่หลากหลายและสอดคล้องกับ CLOs ของรายวิชา รวมถึงสอดคล้องกับวัตถุประสงค์การเรียนการสอน ได้แก่
เช่นการทดสอบปฏิบัติการในห้องปฏิบัติการ (Lab/Practical Test) สำหรับวิชา 4095101 และ 4095102 เพื่อประเมิน CLO ด้านการประยุกต์ใช้ การทำรายงานโครงงาน (Project Report) ในวิชา 4095101 และ 7015908 เพื่อประเมินความสามารถด้านการวิเคราะห์และการสร้างสรรค์ การนำเสนอสัมมนาและการประเมินโดยเพื่อน (Seminar Presentation + Peer Review) ในวิชา 7015907 เพื่อประเมินทักษะการสื่อสารและจริยธรรม การจัดทำข้อเสนอโครงร่างวิจัย (Research Proposal) ในวิชา 7015906 เพื่อประเมินความสามารถด้านระเบียบวิธีวิจัย และการสอบป้องกันวิทยานิพนธ์/สารนิพนธ์ (Thesis/IS Defense) ซึ่งใช้ยืนยันการบรรลุ PLO1 และ PLO3 โดยตรง

% ------------------------------------------------------------------------------
\criterionsub{4.2 The assessment and assessment-appeal policies are shown to be explicit, communicated to students, and applied consistently.}

มหาวิทยาลัยมีนโยบายเรื่องการวัดประเมินผลและการอุทธรณ์ผลการประเมิน ตามประกาศมหาวิทยาลัยราชภัฏอุตรดิตถ์ เรื่อง ระบบและแนวทางการอุทธรณ์เกี่ยวกับผลการประเมินผลการศึกษาของนักศึกษา พ.ศ.~2567 (\url{https://academic.uru.ac.th/document/06943.pdf}) ซึ่งสื่อสารไปยังนักศึกษาผ่านคู่มือนักศึกษาและเว็บกองบริการการศึกษา และนำไปใช้กับนักศึกษาทุกคนอย่างเท่าเทียม

% ------------------------------------------------------------------------------
\criterionsub{4.3 The assessment standards and procedures for student progression and degree completion, are shown to be explicit, communicated to students, and applied consistently.}

มาตรฐานและกระบวนการติดตามความก้าวหน้าและการสำเร็จการศึกษาของนักศึกษาเป็นไปตามข้อบังคับมหาวิทยาลัยราชภัฏอุตรดิตถ์ ว่าด้วยการจัดการศึกษาระดับบัณฑิตศึกษา พ.ศ.~2566 (\url{https://academic.uru.ac.th/document/02898.pdf}) ครอบคลุม
ทั้งเกณฑ์ผ่านรายวิชา (เช่น เกรด B ขึ้นไปสำหรับวิชาบังคับบัณฑิตศึกษา) ระบบสอบประมวลความรู้/การสอบป้องกันวิทยานิพนธ์หรือสารนิพนธ์ และเกณฑ์การตีพิมพ์/การนำเสนอผลงานก่อนสำเร็จการศึกษา (ตามข้อกำหนดของหลักสูตรและประกาศที่เกี่ยวข้อง)

เผยแพร่ผ่านเว็บกองบริการการศึกษาและคู่มือนักศึกษา และใช้บังคับกับทุกคนอย่างสม่ำเสมอ

% ------------------------------------------------------------------------------
\criterionsub{4.4 The assessments methods are shown to include rubrics, marking schemes, timelines, and regulations, and these are shown to ensure validity, reliability, and fairness in assessment.}

หลักสูตรกำหนดให้ทุกรายวิชาต้องมี
\textbf{Rubric} ที่ระบุเกณฑ์การให้คะแนนของแต่ละ CLO \textbf{Marking scheme} ที่แบ่งสัดส่วนคะแนนระหว่างสอบ/ปฏิบัติ/โครงงาน/การมีส่วนร่วม \textbf{Timeline} ที่ระบุวันส่งงานและกำหนดการสอบกลางภาค/ปลายภาคอย่างชัดเจนใน มคอ.3 และ \textbf{Regulations} ที่อ้างอิงข้อบังคับมหาวิทยาลัยฯ และระเบียบการสอบของบัณฑิตศึกษา เพื่อให้การประเมินมีความเที่ยงตรง น่าเชื่อถือ และเป็นธรรม

\marknote{ต้องการหลักฐาน Rubric ฉบับเต็มของรายวิชาบังคับครบทั้ง 6 วิชา ขณะนี้อยู่ระหว่างการจัดทำ}

% ------------------------------------------------------------------------------
\criterionsub{4.5 The assessment methods are shown to measure the achievement of the expected learning outcomes of the programme and its courses.}

วิธีการประเมินในแต่ละรายวิชาออกแบบให้สามารถวัดการบรรลุ CLO และสามารถสะท้อนการบรรลุ PLO ได้ผ่าน CLO--PLO Mapping นอกจากนี้ การสอบป้องกันวิทยานิพนธ์/สารนิพนธ์ใช้คณะกรรมการ 3 คน (โดยมีกรรมการภายนอก 1 คน) ซึ่งประเมินความบรรลุ PLO1 (Create) และ PLO3 (Create Knowledge) โดยตรง

% ------------------------------------------------------------------------------
\criterionsub{4.6 . Feedback of student assessment is shown to be provided in a timely manner.}

หลักสูตรกำหนดช่วงเวลาที่ป้อนกลับผลการประเมินให้ผู้เรียนทราบ
โดยงาน/แบบฝึกหัดจะส่งคืนพร้อมความเห็นภายใน 2 สัปดาห์ การสอบกลางภาคประกาศผลภายใน 2 สัปดาห์พร้อมเปิดโอกาสให้ผู้เรียนตรวจข้อสอบ การสอบปลายภาคประกาศผลผ่านระบบ reg ของมหาวิทยาลัยตามกำหนดการ และความก้าวหน้าวิทยานิพนธ์/สารนิพนธ์ติดตามผ่านการพบนัดอาจารย์ที่ปรึกษาอย่างน้อย 2 ครั้งต่อเดือน

จุดมุ่งหมายคือช่วยเหลือผู้เรียนให้สามารถวางแผนการเรียนทันก่อนสอบตก/drop out

% ------------------------------------------------------------------------------
\criterionsub{4.7 The student assessment and its processes are shown to be continuously reviewed and improved to ensure their relevance to the needs of industry and alignment to the expected learning outcome.}

มหาวิทยาลัยมีระบบทบทวนเกณฑ์การวัดผลตามประกาศมหาวิทยาลัยราชภัฏอุตรดิตถ์ เรื่อง แนวปฏิบัติและการกำกับติดตามการส่งผลการประเมินผลการศึกษาของนักศึกษา พ.ศ.~2567 (\url{https://academic.uru.ac.th/document/06941.pdf}) ดำเนินการในระดับรายวิชาทุกภาค และระดับหลักสูตรประจำปี

\begin{selfassessmentbox}{2}{หลักสูตรมีระบบและประกาศการประเมิน/อุทธรณ์ที่ชัดเจน วิธีประเมินหลากหลายและสอดคล้องกับ CLO/PLO มีระบบป้อนกลับที่ทันเวลา แต่ Rubric ฉบับเต็มของรายวิชาและผลการปรับปรุงจาก feedback ผู้เรียนยังต้องสะสมให้ครบในรอบถัดไป (มีนักศึกษา 4 คน)}
\end{selfassessmentbox}

% ------------------------------------------------------------------------------
\subsection*{รายการหลักฐาน}

หลักฐานประกอบ Criterion~4 ประกอบด้วย AUNQA-4-1 ประกาศมหาวิทยาลัยฯ เรื่อง ระบบและแนวทางการอุทธรณ์ผลการประเมินฯ พ.ศ.~2567, AUNQA-4-2 ข้อบังคับมหาวิทยาลัยฯ ว่าด้วยการจัดการศึกษาระดับบัณฑิตศึกษา พ.ศ.~2566 และ AUNQA-4-3 ประกาศมหาวิทยาลัยฯ เรื่อง แนวปฏิบัติและการกำกับติดตามการส่งผลการประเมินผลการศึกษา พ.ศ.~2567
